% =====================================================================
\section{Introduction}\label{sec:intro}
% =====================================================================
Diabetic retinopathy (DR) is a microvascular complication of diabetes mellitus and one of the leading causes of
preventable vision loss in working-age adults. Pooled population studies estimate that roughly one in three people
with diabetes has some degree of DR and about one in ten has a vision-threatening form \citep{yau2012global}; a more
recent meta-analysis puts the global prevalence among people with diabetes at about 22\% and projects that the
absolute number of affected people will keep rising through 2045 \citep{teo2021global}. Because early DR is
asymptomatic and timely laser or anti-VEGF treatment prevents most severe vision loss, international guidelines
recommend regular screening of every person with diabetes \citep{ting2016diabetic}.

Screening is usually performed with colour fundus photography. A trained grader assigns each image to one of the five
levels of the International Clinical Diabetic Retinopathy (ICDR) severity scale \citep{wilkinson2003proposed}:
no apparent retinopathy (grade 0), mild non-proliferative DR with microaneurysms only (grade 1), moderate
non-proliferative DR (grade 2), severe non-proliferative DR (grade 3) and proliferative DR (grade 4). Grades 2 and
above are generally considered \emph{referable} and require specialist review. Manual grading is accurate when done by
experienced readers, but it is slow, requires a scarce workforce, and shows non-negligible inter-grader variability
\citep{krause2018grader}. In many low- and middle-income regions, where the diabetic population is growing fastest,
the number of trained graders is far below what population screening would require \citep{ting2016diabetic}.

Deep learning has changed what is technically possible. Convolutional neural networks trained on tens of thousands
of graded photographs reach sensitivity and specificity comparable to ophthalmologists for referable DR
\citep{gulshan2016development,gargeya2017automated,ting2017development}, systems trained on very large cohorts cover
the whole disease spectrum \citep{dai2021deep}, and an autonomous AI system has been evaluated in a prospective
pivotal trial in primary care \citep{abramoff2018pivotal}. A systematic review found the diagnostic performance of
deep learning to be broadly equivalent to that of health-care professionals, while warning that few studies used
external validation or reported their methods completely \citep{liu2019comparison}.

\subsection*{Why grading on small datasets remains hard}
These successes rely on large, often private datasets. Many research groups and most new clinical sites have access
only to a few hundred labelled images, such as the Indian Diabetic Retinopathy Image Dataset (IDRiD)
\citep{porwal2018idrid,porwal2020idrid}. In this regime three properties of the task become limiting.

\textbf{(i) Small, sparse evidence.} Adjacent ICDR grades are separated by small findings. A single microaneurysm,
a few pixels wide in a downsampled image, moves an image from grade 0 to grade 1; the number and distribution of
haemorrhages and venous beading separate grade 2 from grade 3; and neovascularisation defines grade 4. A network that
summarises the whole image by global average pooling gives every location the same weight, so a handful of
lesion pixels can be diluted by the large area of healthy retina \citep{wang2017zoom,sun2021lesion}.

\textbf{(ii) Ordered labels.} The grades are ordinal. Predicting grade 4 for a healthy eye is clinically far worse
than predicting grade 1, but standard cross-entropy penalises both errors equally. The quadratic weighted kappa
(QWK) used by most benchmarks \citep{cohen1968weighted} reflects this ordering, and so should the model and its
decision rule \citep{niu2016ordinal,cao2020rank,delatorre2018weighted}.

\textbf{(iii) Imbalance and instability.} Mild and severe grades are under-represented, and with few images per class
both the learned weights and any post-hoc decision thresholds are noisy. Small datasets also make it tempting,
and easy, to tune choices on the test set, which inflates reported performance
\citep{varma2006bias,kapoor2023leakage}.

\subsection*{Our approach}
We address these three issues jointly in a compact model, NITSCI, and in an evaluation protocol designed for small
datasets. For (i), the backbone features are processed by two experts: a \emph{global expert} that relates all
retinal regions through self-attention, and a \emph{lesion expert} that sees the same features re-weighted by a
learned lesion-attention map. Rather than fusing them by simple concatenation, the experts exchange information
through bidirectional cross-attention, so that lesion evidence is interpreted in its anatomical context and global
context is focused on the lesions. An image-adaptive gate, in the spirit of mixtures of experts
\citep{jacobs1991adaptive,shazeer2017outrageously}, decides how much each expert contributes to the final
representation. For (ii), a classification head and an ordinal regression head are trained together and combined
into one continuous grade score. For (iii), the model is first trained on the larger APTOS-2019 dataset
\citep{aptos2019} and then fine-tuned on IDRiD with five-fold cross-validation; class-balanced loss weights, weight
averaging and a fold ensemble with test-time augmentation stabilise the predictions; and every data-dependent choice,
including the decision thresholds, is made on out-of-fold predictions while a held-out test set is used exactly once.

\paragraph{Contributions.} The main contributions of this work are:
\begin{enumerate}
\item \textbf{A cross-expert lesion-aware head.} We propose a head that splits backbone tokens into a global and a
  lesion stream, couples them through bidirectional cross-attention, and fuses them with an image-adaptive gate. The
  lesion stream is learned from grade labels alone, without lesion annotations. The head adds 1.40\,M parameters to
  a 19.85\,M-parameter EfficientNetV2-S backbone.
\item \textbf{An ordinal dual-head decision rule.} The expected grade of the classification head and the output of a
  regression head are averaged into a continuous score whose four cut-points are fitted on out-of-fold data.
\item \textbf{A leakage-free protocol for small datasets.} Cross-dataset pretraining, stratified five-fold
  fine-tuning, an ensemble with test-time augmentation, and a held-out test set that is never used for training,
  model selection or threshold fitting; results are reported with bootstrap confidence intervals and paired McNemar
  tests.
\item \textbf{Component-level evidence.} Seven controlled ablations, run with identical splits and seeds, isolate the
  effect of the lesion map, cross-attention, gate, ordinal head, pretraining and pre-processing on both five-class
  grading and referable-DR screening.
\end{enumerate}

The remainder of the paper is organised as follows. Section~\ref{sec:related} reviews related work.
Section~\ref{sec:data} describes the datasets, data partitioning and pre-processing. Section~\ref{sec:method}
presents the NITSCI architecture, training objective and decision rule. Section~\ref{sec:setup} gives implementation
details, metrics and statistical methods. Section~\ref{sec:results} reports the results and ablations,
Section~\ref{sec:discussion} discusses their implications and limitations, and Section~\ref{sec:conclusion} concludes.
\sectionend{intro}
